\documentclass[../main]{subfiles}
\begin{document}


\section{Codificación de equipos}

Una vez elaborada la lista de equipos es muy importante identificar cada
uno de los equipos con un código único. Esto facilita su localización, su referencia en órdenes de trabajo, en planos, permite la elaboración de registros históricos de fallos e intervenciones, permite el cálculo de indicadores referidos
a áreas, equipos, sistemas, elementos, etc., y permite el control de costes.
Básicamente, existen dos posibilidades a la hora de codificar:
\begin{enumerate}
\item Sistemas de codificación no significativos: son sistemas que asignan un número o un código correlativo a cada equipo, pero el número o código no aporta ninguna información adicional.
    \item  Sistemas de codificación significativos o inteligentes, en el que el código asignado aporta información.
\end{enumerate}

La ventaja del empleo de un sistema de codificación no significativo, de
tipo correlativo, es la simplicidad y la brevedad del código. Con apenas 4 dígitos es posible codificar la mayoría de las plantas industriales. La desventaja es la dificultad para ubicar una máquina a partir de su código: es necesario
tener siempre a mano una lista para poder relacionar cada equipo con su código. Eso, o tener una memoria prodigiosa.
Un sistema de codificación significativo aporta valiosa información sobre
el equipo al que nos referimos: tipo de equipo, área en el que está ubicada, familia a la que pertenece, y toda aquella información adicional que queramos
incorporar al código. El problema es que al añadir más información el código
aumenta de tamaño.
Como quiera el empleo de sistemas correlativos es muy sencillo, estudiaremos los sistemas de codificación significativos.
Información útil que debe contener el código de un ítem
La información que debería contener el código de un equipo debería ser
el siguiente:
\begin{itemize}
\item Planta a la que pertenece.
\item Area al que pertenece dentro de la planta.
\item Tipo de equipo.
Los elementos que forman parte de un equipo deben contener información adicional:
\item Tipo de elemento.
\item Equipo al que pertenecen.
\item Dentro de ese equipo, sistema en el que están incluidos.
\item  Familia a la que pertenece el elemento. La clasificación en familias es
muy útil, ya que nos permite hacer listados de elementos. Se puede encontrar una lista de familias en que pueden clasificarse los elementos más adelante.
\end{itemize}
Una vez elaborada la lista de equipos, y teniendo en cuenta todas las consideraciones anteriores, es posible abordar la tarea de la codificación, fijando
los criterios que la regirán. Un ejemplo de codificación se incluye en el punto siguiente.
\ejemplo{ Una forma de codificar los activos de una planta}{
Para este ejemplo, se ha utilizado la planta de producción de Compact
Disc que se detallaba en el apartado dedicado a la elaboración de listas de
equipos. Se han codificado tanto los equipos como los elementos que los
componen.}
Para definir el código en este ejemplo, se ha definido la siguiente estructura:

\begin{figure}[ht]
    \centering
    \includegraphics[width=0.3\textwidth]{images/04/codificacionimage.png}
    \caption{Códigos para equipos}
\end{figure}

Como se indica en la Figura, el Área de la Planta en que está ubicado
el equipo estaría definido por dos caracteres alfanuméricos, el tipo de
equipo por dos caracteres alfabéticos, y el número correlativo por dos caracteres numéricos.

En la siguiente Tabla figuran algunos ejemplos:

\begin{table}[h!]
    \centering
    \begin{tabular}{|c|l|c|l|}\hline
Código & Tipo de equipo & Código & Tipo de equipo\\\hline
BS & Caldera de vapor & C0 &Red de Aire comprimido \\
BF & Caldera de Fluido Térmico & R0 & Red de Agua de refrigeración \\
CR & Compresor rotativo &F0 &Red contraincendios\\
CC & Compresor Centrífugo & E0 & Red eléctrica general\\
CT & Compresor de tornillo & V0 & Red de Vacío\\
SC &Secador de aire por sistema frig. & S0 & Red de Vapor \\
SG &Secador de aire por silica gel &H0& Red de agua sobrecalentada \\
TS &Tanque de alm. a presión &W0& Red de agua industrial
atm.\\ \hline


    \end{tabular}
    \caption{Ejemplos de códigos que pueden
utilizarse para identificar el tipo de equipo}
    
\end{table}

\begin{table}[h!]
    \centering
    \begin{tabular}{|c|l|c|l|}\hline
Código & Tipo de equipo & Código & Tipo de equipo\\\hline
TP &Tanque de alm. a presión &G0& Red de gas
superior a la atm.\\
OE& Horno eléctrico &FC& Equipo frig. por compresión\\
OG& Horno de gas &FA& Equipo frig. por absorción\\
TR& Reactor &CL& Climatizador\\
ET& Intercambiador de haces tubulares &EP& Intercambiador de placas\\
AR& Aerorefrigerador &IN& Inyectora\\
RI& Torre Refrig. tiro inducido &RN& T. Refrig. tiro natural\\
RF& Torre Refrig. tiro natural &TG& Turbina de gas\\
TV& Turbina de Vapor &PS& Impresora por serigrafía\\
PM& Envasadora & & \\ \hline
    \end{tabular}
    \caption{Ejemplos de códigos que pueden
utilizarse para identificar el tipo de equipo (continuación)}
   
\end{table}

Utilizando esta Tabla de tipos de equipo, y teniendo en cuenta la lista
de equipos que contenía la planta del apartado anterior, la codificación resultante sería la de el la tabla.

\begin{table}[ht]
    \centering
    \begin{tabular}{|c|c|c|c|c|c|l|}\hline
C&ó&d&i&g&o&Descripción\\\hline
1&1&F&1& & &Equipo contraincendios \\
1&1&F&0& & &Red de agua contraincendios\\
1&1&C&T&0&1&Compresor 1\\
1&1&C&T&0&2&Compresor 2\\
1&1&S&C&0&1&Secador de aire\\
1&1&C&0& & &Red de aire comprimido\\
1&1&V&C&0&1&Bomba de Vacío\\
1&1&V&0& & &Red de Vacío\\
1&2&L&P&0&1&Laboratorio de Premastering\\
1&2&L&M&0&1&Laboratorio de Mastering\\
1&2&L&T&0&1&Laboratorio de Matrices\\
1&3&I&N&0&1&Inyectora 1\\
1&3&I&N&0&2&Inyectora 2\\
1&3&I&N&0&3&Inyectora 3\\
1&4&P&S&0&1&Impresora 1\\
1&4&P&S&0&2&Impresora 2\\
1&4&L&S&0&1&Laboratorio de serigrafía\\
1&5&P&M&0&1&Envasadora 1\\
1&5&P&M&0&2&Envasadora 2\\
\hline
    \end{tabular}
    \caption{Codificación resultante.}

\end{table}

\clearpage
\newpage

\subsubsection*{Códigos para elementos}
El código de un elemento que forma parte de un equipo estaría formado en este ejemplo por un total de 17 caracteres, con la siguiente estructura:
\begin{itemize}
\item Los 6 primeros identificarían el equipo, tal y como se ha detallado anteriormente.
\item Un carácter más alfabético identificaría la familia a la que pertenece el elemento.
\item Los tres caracteres siguientes identificarían el sistema.
\item Los caracteres siguientes, hasta 7 (longitud variable), serían caracteres alfanuméricos, que identificarían las características del
elemento y aportarían un número correlativo.
\item  Un ultimo carácter, de aplicación exclusiva para el caso de redundancia (elementos duplicados, triplicados, etc.).
En la siguiente Figura puede verse la estructura del código de un elemento:
\begin{figure}[h]
    \centering
    \includegraphics{images/04/codificacionelementos.png}
    \caption{Código de elementos}
\end{figure}

\end{itemize}


Las familias a las que puede pertenecer un elemento pueden ser las siguientes:

\begin{table}[h]
    
    \centering 
    
    \begin{tabular}{|c|l|}\hline
Código&Familia\\\hline
B&Bomba\\
M&Motor\\
V&Válvula\\
I&Instrumento\\
C&Componente de cuadro eléctrico\\
E&Elemento eléctrico\\
P&Pieza mecánica\\
T&Tubería\\
F&Filtro\\
N&Cilindros y actuadores neumáticos (no válvulas)\\
H&Cilindros y actuadores hidráulicos\\
O&Brida\\\hline
    \end{tabular}
    \caption{Familias de elementos}
\end{table}


Indicar la familia a la que pertenece el elemento tiene una ventaja, nos
permite hacer listados de válvulas, motores, bombas, instrumentos, etc. Estas listas pueden ser interesantes en muchos casos; por ejemplo, si quisiéramos saber cuántos y qué motores tenemos en la planta, para estudiar su
posible estandarización. O saber cuántos instrumentos de medida hay instalados, para hacer un Plan de Calibración.

\actividad{Contestar, investigar y responder.}
\begin{enumerate}
    \item ¿Cuál es la diferencia entre un sistema de codificación significativo y no significativo?
    \item ¿Cuáles son las ventajas y desventajas para el sistema de codificación significativo y no significativo.?
    \item Un sistema de codificación no significativo es el código QR y código de Barras. Investigar y explicar como es el funcionamiento de ambos sistemas.
    \begin{enumerate}
        \item ¿Para qué está la norma de código de barras EAN13?
        \item ¿Qué es un código bidimensional?
    \end{enumerate}
    
    \item Elaborar el diseño de un \emph{código significativo} de equipo para una máquina del taller de la escuela. 
    
\end{enumerate}


\end{document}
